%! Conditional Probability, Independence, and Simulation — Unit 3, Lesson 3.4 — Student Packet Cover Sheet
% THREE packet rows — Warm-Up, Guided Notes & Practice, Homework. A contents list, no Score column.
\documentclass[12pt]{article}
\usepackage{saar-article}
\usepackage{saar-boxes}
\usepackage{ltablex}
\keepXColumns
% Measured banner: the band grows to fit the title block, so a title long enough to wrap grows
% the band rather than running out of it (the stats course's \coverbanner; saar-article does not
% carry one, so it is defined here — every cover copies this block verbatim).
\newsavebox{\bannerbox}
\newlength{\bannerht}
\newlength{\coverbannerpad}
\setlength{\coverbannerpad}{0.16in}       % breathing room above and below the text
\newcommand{\coverbanner}[2]{%
  \sbox{\bannerbox}{%
    \begin{minipage}{\dimexpr\paperwidth-1.4in\relax}%
      \centering
      {\color{white}\textbf{\LARGE \CourseName}}\\[4pt]
      {\color{frostmid}\large\textbf{#1}}\\[6pt]
      {\color{white}\textbf{\large #2}}%
    \end{minipage}}%
  \setlength{\bannerht}{\dimexpr\ht\bannerbox+\dp\bannerbox+2\coverbannerpad\relax}%
  \noindent\begin{tikzpicture}[remember picture, overlay]
    \fill[cerulean] (current page.north west)
      rectangle ([yshift=-\bannerht]current page.north east);
    \node[anchor=north, inner sep=0pt]
      at ([yshift=-\coverbannerpad]current page.north) {\usebox{\bannerbox}};
  \end{tikzpicture}\par
  % The text body starts 0.6in below the page edge (geometry top=0.6in), so reserve only what
  % the band overruns that, plus a gap under the band.
  \vspace*{\dimexpr\bannerht-0.6in+0.16in\relax}%
}

\begin{document}

% Full-bleed cerulean banner, sized to the title block.
\coverbanner{Unit 3}{Lesson 3.4 \quad Conditional Probability, Independence, and Simulation}

\vspace{0.06in}
% NAMESTRIP: this is the ONLY component that carries the name row.
\namedateperiod
\vspace{0.18in}

% The standard codes stay in the teacher-facing plan.
\begin{learningtargetbox}
\small
\begin{enumerate}[leftmargin=1.5em, itemsep=5pt]
  \item \ldots{} find a \textbf{conditional probability} from a two-way table by restricting to
        the ``given'' row or column.
  \item \ldots{} use $P(B \mid A) = \dfrac{P(A \text{ and } B)}{P(A)}$ with a Venn diagram, and
        know it is $0$ for mutually exclusive events.
  \item \ldots{} test whether two events are \textbf{independent} by comparing $P(A \mid B)$ with
        $P(A)$.
  \item \ldots{} read a \textbf{simulation}, explain the \textbf{Law of Large Numbers}, and justify a
        decision with the conditional probability that answers the question.
\end{enumerate}
\end{learningtargetbox}

\vspace{0.14in}

\begin{tocbox}
\small
\renewcommand{\arraystretch}{1.5}
\begin{tabularx}{\linewidth}{@{} c l X @{}}
\rowcolor{frost}
\textbf{\#} & \textbf{Component} & \textbf{Description} \\
\midrule
1 & Warm-Up & A two-way table from a summer coding course: reading counts, percent of a part,
             comparing decimals \\
2 & Guided Notes \& Practice & Conditional probability from a table and by formula; the
             independence test; a simulated study app and a decision --- then a spam filter we
             work together \\
3 & Homework & A free-throw simulation and a coach's decision. \textbf{Scored 1 to 5, assigned
             today, due the first class after two study halls.} \\
\end{tabularx}
\end{tocbox}

\vspace{0.14in}

% Keep in Mind carries the lesson's CONTENT — the rule, the test, the trap — never its process.
\begin{remindbox}
\small
The word after the bar names the group: in $P(A \mid B)$, restrict to $B$ and divide by $B$'s total.
Events are \textbf{independent} when $P(A \mid B) = P(A)$; if the two differ, the events are
dependent. By the \textbf{Law of Large Numbers}, a long \textbf{simulation} settles near the actual
probability. \textbf{Watch out:} $P(A \mid B)$ and $P(B \mid A)$ are different numbers. Before you
decide, name the group the question is about, and put that group after the bar.
\end{remindbox}

\end{document}
